\begin{center}[6.]\par\medskip\textit{Zur mathematischen Theorie der elektrodynamischen Wirkungen.}\end{center}
1.\quad Die gegenseitige Wirkung zweier Stromelemente $\mathrm ds,\mathrm ds'$ aufeinander, wobei $i,i'$ die Intensitäten der Ströme bezeichnen, drückt \textsc{Ampère} durch die Formel
\[
ii'(\sin\theta\sin\theta'\cos\omega+k\cos\theta\cos\theta')r^n\mathrm ds\,.\,\mathrm ds'
\]
aus. Dabei setzt er voraus, dass sie entlang der verbindenden Geraden stattfindet. Positive Zeichen stehen für Anziehung, negative für Abstoßung. Hier bezeichnet $r$ den Abstand der Elemente voneinander, $\theta$ und $\theta'$ die Winkel der Elemente $\mathrm ds,\mathrm ds'$ mit $r$, wobei diese Linie in beiden Fällen in derselben Richtung genommen wird, und schließlich $\omega$ den Winkel zwischen den Ebenen durch $r$ und $\mathrm ds$ einerseits und $r$ und $\mathrm ds'$ andererseits. Aus seinen Versuchen schloss \textsc{Ampère}, dass man $n=-2$, $k=-\frac{1}{2}$ setzen muss. Die gegenseitige Anziehung wird also durch
\[
\frac{ii'(\sin\theta\sin\theta'\cos\omega-\frac{1}{2}\cos\theta\cos\theta')\mathrm ds\,.\,\mathrm ds'}{rr}
\]
gemessen. Ein negativer Wert dieses Ausdrucks bedeutet Abstoßung.

2.\quad Wir bezeichnen außerdem mit $a,a'$ die gesamten Stromlängen, mit $\rho$ die Quadratwurzel aus $r$, mit $x,y,z$ die Koordinaten eines beliebigen Punkts in $s$, mit $x',y',z'$ die Koordinaten eines beliebigen Punkts in $s'$ und mit $\varepsilon$ den
\SourceLeaf{613}Winkel, den zwei Stromelemente $\mathrm ds,\mathrm ds'$ miteinander bilden. Partielle Differentiationen nach $s$ und $s'$ unterscheiden wir durch die Zeichen $\mathrm d,\mathrm d'$, partielle Integrationen entsprechend durch die Zeichen $\int,\int'$. Wir haben
\[
\begin{gathered}
\mathrm dr=2\rho\,\mathrm d\rho=\cos\theta\,.\,\mathrm ds,\qquad \mathrm d'r=2\rho\,\mathrm d'\rho=-\cos\theta'\,.\,\mathrm ds'\\
\rho^4=(x'-x)^2+(y'-y)^2+(z'-z)^2\\
2\rho^3\mathrm d\rho=-(x'-x)\mathrm dx-(y'-y)\mathrm dy-(z'-z)\mathrm dz\\
\begin{aligned}
6\rho\rho\,\mathrm d\rho\,.\,\mathrm d'\rho+2\rho^3\mathrm d\mathrm d'\rho
&=-\mathrm dx\,.\,\mathrm d'x'-\mathrm dy\,.\,\mathrm d'y'-\mathrm dz\,.\,\mathrm d'z'\\
&=-\cos\varepsilon\,.\,\mathrm ds\,.\,\mathrm ds'\\
2\rho^3\mathrm d\mathrm d'\rho&=(-\cos\varepsilon+\tfrac{3}{2}\cos\theta\,.\,\cos\theta')\mathrm ds\,.\,\mathrm ds'\\
&=-(\sin\theta\sin\theta'\,.\,\cos\omega-\tfrac{1}{2}\cos\theta\,.\,\cos\theta')\mathrm ds\,.\,\mathrm ds'
\end{aligned}
\end{gathered}
\]
Der obige \textsc{Ampère}sche Ausdruck nimmt also folgende Form an
\[
-\frac{2ii'\mathrm d\mathrm d'\rho}{\rho}
\]

3.\quad Mit dem vorangestellten Zeichen $\delta$ bezeichnen wir die unendlich kleinen Variationen einer Größe, soweit diese von unendlich kleinen virtuellen Ortsänderungen der Punkte der Leiter $s,s'$ abhängen. Virtuelle Ortsänderungen sind alle willkürlich gedachten Ortsänderungen, die mit den Bedingungen vereinbar sind, die sich aus der Beschaffenheit der Leiter und ihren äußeren Verhältnissen ergeben.

4.\quad Die gesamte Wirkung der Ströme aufeinander ist als bekannt anzusehen, wenn man andere Kräfte angeben kann, die auf endlich oder unendlich viele einzelne Punkte der Ströme wirken und ihnen in dem Sinne äquivalent sind, dass die entgegengesetzten Kräfte ihnen das Gleichgewicht halten. Sei $W$ die Summe dieser letzteren Kräfte, jeweils mit ihren virtuellen Bewegungen multipliziert. Nach dem Prinzip der virtuellen Bewegungen muss dann
\[
W+\int\int'\delta r\,.\,\frac{2ii'\mathrm d\mathrm d'\rho}{\rho}=W+4ii'\int\int'\mathrm d\mathrm d'\rho\,.\,\delta\rho
\]
$=0$ sein oder wenigstens keinen positiven Wert annehmen können, und zwar für jedes System virtueller Bewegungen der Leiter, das mit den gegebenen Bedingungen vereinbar ist. Für sich allein halten die Stromkräfte die Leiter dagegen im Gleichgewicht, wenn $\int\int'\delta\rho\frac{\mathrm d\mathrm d'\rho}{\rho}$ für alle virtuellen Bewegungen $=0$ oder negativ ist. Die Integrationen sind hier von $s=0$ bis $s=a$ und von $s'=0$ bis $s'=a'$ auszuführen.

\SourceLeaf{614}5.\quad Um die Natur dieses Integrals zu verstehen, entwickeln wir die Variation von $\int\int'\mathrm d\rho\,.\,\mathrm d'\rho$. Wir haben
\[
\begin{aligned}
\delta\int\int'\mathrm d\rho\,.\,\mathrm d'\rho&=\int\int'\delta(\mathrm d\rho\,.\,\mathrm d'\rho)\\
&=\int\int'\delta\mathrm d\rho\,.\,\mathrm d'\rho+\int\int'\delta\mathrm d'\rho\,.\,\mathrm d\rho
\end{aligned}
\]
Nun gilt aber
\[
\mathrm d\int'\mathrm d'\rho\delta\rho=\int'\mathrm d\delta\rho\,.\,\mathrm d'\rho+\int'\delta\rho\,.\,\mathrm d\mathrm d'\rho
\]
oder, wenn man von $s=0$ bis $s=a$ integriert und die Werte von $\int'\mathrm d'\rho\,.\,\delta\rho$ für $s=0$ und $s=a$ mit $F',G'$ bezeichnet,
\[
-F'+G'=\int\int'\mathrm d\delta\rho\,.\,\mathrm d'\rho+\int\int'\delta\rho\,.\,\mathrm d\mathrm d'\rho
\]
und ebenso, wenn man die Werte von $\int\mathrm d\rho\,.\,\delta\rho$ für $s'=0$ und $s'=a'$ mit $F,G$ bezeichnet,
\[
-F+G=\int\int'\mathrm d'\delta\rho\,.\,\mathrm d\rho+\int\int'\delta\rho\,.\,\mathrm d\mathrm d'\rho
\]
Folglich
\[
\begin{aligned}
-F-F'+G+G'&=\int\int'\mathrm d\delta\rho\,.\,\mathrm d'\rho+\int\int'\mathrm d'\delta\rho\,.\,\mathrm d\rho\\
&\qquad+2\int\int'\delta\rho\,.\,\mathrm d\mathrm d'\rho\\
&=\delta\int\int'\mathrm d\rho\,.\,\mathrm d'\rho+2\int\int'\delta\rho\,.\,\mathrm d\mathrm d'\rho
\end{aligned}
\]

Das virtuelle Moment der gegenseitigen Einwirkung der Stromabschnitte $s,s'$ aufeinander ist also
\[
\begin{aligned}
V=4ii'\int\int'\mathrm d\mathrm d'\rho\,.\,\delta\rho&=-2ii'\delta\int\int'\mathrm d\rho\,.\,\mathrm d'\rho\\
&\qquad+2ii'(G+G'-F-F')
\end{aligned}
\]

6.\quad Ein Strom übt eine stromerzeugende oder elektromotorische Kraft nur aus, wenn er entsteht oder sich bewegt. Die elektromotorische Kraft eines Stromelements $i\,\mathrm ds$ wirkt in jedem Punkt mit einer Stärke, die zur Entfernung $r$ umgekehrt und zum auf diese Linie projizierten Stromelement direkt proportional ist. Sie wirkt entlang der Linie $r$ selbst, aber stets in entgegengesetzter Richtung. Daraus folgt leicht, dass sich das virtuelle Moment der elektromotorischen Kraft des Elements $i\,\mathrm ds$ auf zwei Arten ausdrücken lässt.
\[
+\frac{i\,\mathrm ds}{r}\,.\,\frac{\mathrm dr}{\mathrm ds}\,.\,\delta r
\]
oder auch
\SourceLeaf{615}\[
4i\,\mathrm d\rho\,.\,\delta\rho
\]

Im Leiterelement $\mathrm ds'$ ist die elektromotorische Kraft des Stromelements $i\,\mathrm ds$ also
\[
=4i\,\mathrm d\rho\,.\,\frac{\mathrm d\rho}{\mathrm ds'}\,.\,\mathrm ds'=4i\,\mathrm d\rho\,.\,\mathrm d'\rho
\]

Das Doppelintegral $4i\int\int'\mathrm d\rho\,.\,\mathrm d'\rho$ ist also die gesamte elektromotorische Kraft im Leiterabschnitt von $s'=0$ bis $s'=a'$, die der Stromabschnitt $is$ von $s=0$ bis $s=a$ ausübt. Wir bezeichnen sie mit $A$. Ferner ist das virtuelle Moment der elektromotorischen Kraft, die derselbe Stromabschnitt an einem gegebenen Punkt ausübt,
\[
=4i\int\mathrm d\rho\,.\,\delta\rho\qquad\text{von }s=0\text{ bis }s=a
\]
Ist dieses am Anfang des Leiterabschnitts $s'$ $=B$ und an dessen Ende $=C$, so gilt
\[
B=4iF,\qquad C=4iG
\]
Bezeichnen ebenso $B',C'$ die virtuellen Momente der elektromotorischen Kraft, die ein Stromabschnitt $i's'$ von $s'=0$ bis $s'=a'$ am Anfangs- und am Endpunkt von $s$ ausübt, so ist $B'=4i'F'$, $C'=4i'G'$.

Wir erhalten folglich
\[
V=-\tfrac{1}{2}i'\delta A-\tfrac{1}{2}i'B+\tfrac{1}{2}i'C-\tfrac{1}{2}iB'+\tfrac{1}{2}iC'
\]

